\documentclass[11pt,a4paper]{article}
\usepackage{harmonikabau_preamble}

\newcommand{\hbdoknr}{0004}
\newcommand{\hbtitel}{Reed Frequency Variation}
\newcommand{\hbuntertitel}{through Chamber Coupling -- Two-Filter Model}
\newcommand{\hbheadtext}{Doc.\,0004 --- Frequency Variation: Two-Filter Model}
\newcommand{\hbfoottext}{Cross-references: Doc.\,0002 $\cdot$ Doc.\,0003}

\AtBeginDocument{%
  \fancyhead[L]{\small\hbheadtext}%
  \fancyhead[R]{\small\thepage}%
  \fancyfoot[C]{\small\color{hb@darkgray}\hbfoottext}%
}

\begin{document}
\hbmaketitle

{\small\itshape Shows why the chamber shifts the reed pitch despite separate natural frequencies, and how this can be calculated as a filter combination of two coupled systems. Numerical values: 50\,Hz bass reed, $f_H = 274$\,Hz, $Q_H \approx 7$ (from Doc.\,0002).}

\vspace{2mm}
\begin{center}\small
\begin{tabularx}{\textwidth}{l X}
\hbtoprule
\textbf{Document References} \\
\midrule
Doc.\,0001 & Analysis of instrument acoustics and psychoacoustics \\
Doc.\,0002 & Flow analysis of bass reed 50\,Hz -- v8 \\
Doc.\,0003 & Impedance comparison: free-beating reed vs.\ labial pipe \\
\textbf{Doc.\,0004} & \textbf{Frequency variation -- two-filter model (this document)} \\
Doc.\,0006 & Symbol reference \\
Doc.\,0007 & Treble reed block -- chamber frequencies \\
Doc.\,0008 & Tone modification by chamber geometry \\
Doc.\,0500 & Guide to aesthetic forensics of accordion casework \\
\midrule
\href{berechnungen_verifikation.py}{berechnungen\_verifikation.py} & Verification script for all calculations \\
\bottomrule
\end{tabularx}
\end{center}

\vspace{4mm}

\section{The Practical Finding}

Every experienced accordion builder knows the phenomenon: \textbf{the same reed, mounted in different chambers, sounds at different pitches.} The frequency variation is small (a few cents) but clearly measurable. Cottingham and Millot measured shifts of 10--30\,cents in free-beating reeds when the resonator was changed. In accordion bass chambers, where $f_\text{reed} \ll f_H$, the effect is smaller (1--5\,cents) but systematic.

This contradicts the simplified statement from Doc.\,0002 Ch.\,11 that the chamber ``does not disturb the fundamental'' (phase $\approx -89°$). The phase is indeed almost $-90°$ --- but \textbf{almost} is not \textbf{exactly}, and even a pure compliance shifts the frequency because it imposes an additional spring stiffness on the reed.

Doc.\,0003 shows that the fundamental difference from the organ pipe lies in the \textbf{time-variant gap} (impedance factor~200 per cycle). Nevertheless, the frequency shift is calculable --- with the following two-filter model.

\begin{warnbox}
\textbf{Practical consequence:} An instrument is not tuned ``on the bench'' and then installed in any chamber. The reed is tuned \textbf{in its chamber}. If the chamber is changed, the pitch changes.
\end{warnbox}

\section{The Two-Filter Model}

\subsection{Two Resonators, One Gap}

\textbf{Filter~1 -- The reed (mechanical band-pass):} Spring-mass system with $f_1 = 50$\,Hz and $Q_1 = 50$--$150$ (Doc.\,0002 Ch.\,9):
\begin{equation}
H_1(f) = \frac{f_1^2}{f_1^2 - f^2 + j f f_1 / Q_1}
\end{equation}

\textbf{Filter~2 -- The chamber (acoustic band-pass):} Helmholtz resonator with $f_H = 274$\,Hz and $Q_H = 5$--$10$ (Doc.\,0002 Ch.\,11):
\begin{equation}
H_2(f) = \frac{f_H^2}{f_H^2 - f^2 + j f f_H / Q_H}
\end{equation}

\textbf{The coupling:} The gap converts reed deflection into volume flow ($q = A_\text{eff} \cdot \dot{x}$) and chamber pressure into force ($F = A_\text{eff} \cdot p$). The effective gap area $A_\text{eff} = W \cdot h_\text{max}/3 = 4$\,mm$^2$ (Doc.\,0002 Ch.\,2) is the coupling element.

\subsection{The Feedback Loop (Barkhausen Condition)}

The system is a \textbf{feedback loop}:

\medskip
\centerline{$x \to S(x) \to q \to p_\text{chamber} \to F_\text{Bernoulli} \to x$}
\medskip

The oscillation frequency satisfies the Barkhausen condition:
\begin{equation}
\varphi_\text{reed}(f) + \varphi_\text{chamber}(f) + \varphi_\text{coupling}(f) = 0
\end{equation}

Without chamber: $\varphi_\text{reed}(f_\text{mech}) = 0$ $\to$ oscillation at $f_\text{mech}$. With chamber: $\varphi_\text{chamber}(f_\text{mech}) \neq 0$ $\to$ the reed shifts its frequency.

\begin{keybox}
\textbf{Core statement:} The chamber adds a phase shift to the feedback. To satisfy the phase condition, the system shifts its oscillation frequency. The shift depends on how steeply the reed phase varies at $f_\text{mech}$ ($= Q_1$) and how large the chamber phase is at $f_\text{mech}$ ($= f_\text{mech}/f_H$ and coupling strength).
\end{keybox}

\section{The Three Mechanisms of Frequency Shift}

\subsection{Mechanism 1: Static Compliance Load (Air Spring)}

The chamber ($V = 180$\,cm$^3$) acts as an additional spring:
\begin{equation}
\Delta k = \frac{A_\text{eff}^2 \cdot \rho c^2}{V} = \frac{(4 \times 10^{-6})^2 \times 1{.}2 \times 343^2}{1{.}8 \times 10^{-4}} = 0{.}013\;\text{N/m}
\end{equation}

Compared with the mechanical spring stiffness $k_\text{mech} = \omega_1^2 m_\text{eff} = 37{.}0$\,N/m (with $m_\text{eff} = 0{.}243 \times m_\text{total} = 0{.}375$\,g for the first cantilever mode):
\begin{equation}
\frac{\Delta f}{f} = \frac{\Delta k}{2 k_\text{mech}} = \frac{0{.}013}{74{.}0} = 1{.}7 \times 10^{-4} \quad \Rightarrow \quad \mathbf{0{.}3\;\text{cents}}
\end{equation}

Always \textbf{upward} (stiffer spring), always present, independent of $f_H$.

\subsection{Mechanism 2: Resonant Amplification at Overtones near $f_H$}

The impedance is amplified by factor $G(f)$ near $f_H$:
\begin{equation}
G(f) = \frac{1}{\sqrt{\left(1 - (f/f_H)^2\right)^2 + \left(f/(f_H \cdot Q_H)\right)^2}}
\end{equation}

\begin{table}[H]
\centering\small
\resizebox{\linewidth}{!}{\begin{tabular}{rrrrl}
\hbtoprule
OT $n$ & $f$ [Hz] & $f/f_H$ & $G(f)$ & Phase $Z$ [°] \\
\midrule
1 & 50 & 0.18 & 1.0 & $-88.5$ \\
3 & 150 & 0.55 & 1.4 & $-83.6$ \\
5 & 250 & 0.91 & 4.7 & $-52.1$ \\
$f_H$ & 274 & 1.00 & 7.0 & $0.0$ \\
6 & 300 & 1.09 & 4.0 & $+51.8$ \\
8 & 400 & 1.46 & 0.9 & $+79.6$ \\
10 & 500 & 1.82 & 0.4 & $+83.6$ \\
\bottomrule
\end{tabular}}
\caption{Impedance amplification $G(f)$ of the Helmholtz resonator at the overtones of the 50\,Hz reed ($f_H = 274$\,Hz, $Q_H = 7$). Phase values from Doc.\,0002 Ch.\,11.}
\end{table}

\subsection{Mechanism 3: Nonlinear Feedback to the Fundamental}

The overtones are bound to the fundamental through nonlinear gap coupling (Doc.\,0002 Ch.\,8, Doc.\,0003 Ch.\,5). If the $n$-th overtone is pulled by $\Delta f_n$:
\begin{equation}
\Delta f_1 \approx \frac{1}{n} \cdot \Delta f_n
\end{equation}

\subsection{The Bernoulli Amplification Factor}

The coupling constant:
\begin{equation}
\kappa = \frac{A_\text{eff}}{\sqrt{m_\text{eff} \cdot C_a \cdot \omega_1 \cdot \omega_2}}
       = \frac{4 \times 10^{-6}}{\sqrt{3{.}75 \times 10^{-4} \cdot 1{.}28 \times 10^{-9} \cdot 314 \cdot 1721}}
\end{equation}
\begin{equation*}
       \approx 0{.}008
\end{equation*}
with acoustic compliance $C_a = V / (\rho c^2) = 1{.}28 \times 10^{-9}$\,m$^3$/Pa.

$\kappa \approx 0{.}008 \ll 1$: \textbf{weak coupling}. The gap area (4\,mm$^2$) is tiny compared to the chamber area ($\sim$950\,mm$^2$). But the Bernoulli mechanism \textbf{amplifies} the coupling. The chamber pressure fluctuation modulates the Bernoulli suction at the gap:
\begin{equation}
\beta = \frac{\delta p_\text{chamber}}{p_\text{Bernoulli}} = \frac{q_\text{reed} \cdot |Z_H(f_1)|}{\frac{1}{2}\rho v_\text{gap}^2} \approx 0{.}63\,\%
\end{equation}

This 0.63\,\% pressure modulation shifts the frequency by:
\begin{equation}
\Delta f_\text{Bernoulli} \approx f_1 \cdot \frac{\beta}{2} = 50 \times 0{.}003 = \mathbf{0{.}16\;\text{Hz} \approx 5{.}4\;\text{cents}}
\end{equation}

\begin{keybox}
Pure acoustic coupling ($\kappa^2$) yields only 0.3\,cents. But Bernoulli amplification ($\beta \approx 0{.}6\,\%$ pressure modulation) raises the frequency shift to \textbf{$\sim$5\,cents} --- clearly audible and consistent with practice. The two filters are not passively cascaded, but \textbf{actively coupled through the Bernoulli mechanism}.
\end{keybox}

\section{The Complete Filter Combination}

The closed transfer function of the total system:
\begin{equation}
\boxed{H_\text{ges}(f) = \frac{H_1(f)}{1 - \kappa_\text{eff} \cdot H_1(f) \cdot H_2(f)}}
\end{equation}

The poles of $H_\text{ges}$ give the actual oscillation frequencies of the coupled 4th-order system:

\textbf{Mode~1 (reed mode):}
\begin{equation}
f_1' \approx f_1 \cdot \left(1 + \kappa^2 \frac{f_1^2}{f_H^2 - f_1^2}\right) = 50{.}000\;\text{Hz}
\end{equation}
$(+0.3$ cents purely acoustic)

\textbf{Mode~2 (chamber mode):}
\begin{equation}
f_2' \approx f_H \cdot \left(1 - \kappa^2 \frac{f_H^2}{f_H^2 - f_1^2}\right) = 273{.}98\;\text{Hz}
\end{equation}

\textbf{Mode repulsion:} The modes repel each other --- the lower is pushed down (after Bernoulli correction: \textit{up}), the upper is pushed up. With Bernoulli amplification ($\beta = 0{.}63\,\%$), the shift of the reed mode becomes \textbf{$\sim$5\,cents}.

\section{Parametric Dependencies}

\subsection{Chamber Volume $V$}

$\Delta k \sim 1/V$: smaller chamber $\to$ stiffer air spring $\to$ larger shift.

\begin{table}[H]
\centering\small
\resizebox{\linewidth}{!}{\begin{tabular}{rrrrr}
\hbtoprule
$V$ [cm$^3$] & $f_H$ [Hz] & $\Delta k$ [N/m] & Static [cents] & With Bernoulli [cents] \\
\midrule
90 & 387 & 0.025 & 0.6 & $\sim$11 \\
120 & 335 & 0.019 & 0.4 & $\sim$8 \\
180 & 274 & 0.013 & 0.3 & $\sim$5 \\
240 & 237 & 0.009 & 0.2 & $\sim$4 \\
360 & 194 & 0.006 & 0.1 & $\sim$3 \\
\bottomrule
\end{tabular}}
\caption{Frequency shift as a function of chamber volume}
\end{table}

\subsection{Gap Area $A_\text{eff}$ (Deflection)}

$\Delta k \sim A_\text{eff}^2$: more deflection $\to$ quadratically stronger coupling.

\begin{table}[H]
\centering\small
\resizebox{\linewidth}{!}{\begin{tabular}{rrrr}
\hbtoprule
$h_\text{max}$ [mm] & $A_\text{eff}$ [mm$^2$] & Static [cents] & $\kappa$ \\
\midrule
0.5 & 1.3 & 0.03 & 0.003 \\
1.0 & 2.7 & 0.13 & 0.005 \\
1.5 & 4.0 & 0.29 & 0.008 \\
2.0 & 5.3 & 0.52 & 0.010 \\
2.5 & 6.7 & 0.81 & 0.013 \\
\bottomrule
\end{tabular}}
\caption{Frequency shift as a function of reed deflection (chamber 180\,cm$^3$)}
\end{table}

\subsection{Helmholtz Frequency $f_H$}

Maximum coupling when $f_H = n \cdot f_\text{reed}$ (overtone exactly on resonance).

\begin{table}[H]
\centering\small
\resizebox{\linewidth}{!}{\begin{tabular}{rlrrl}
\hbtoprule
$f_H$ [Hz] & Nearest OT & Distance [Hz] & $G$ & Effect \\
\midrule
150 & 3rd OT (150) & $\pm 0$ & 7.0 & \textbf{Strong -- critical} \\
200 & 4th OT (200) & $\pm 0$ & 7.0 & \textbf{Strong -- critical} \\
250 & 5th OT (250) & $\pm 0$ & 7.0 & \textbf{Strong -- critical} \\
274 & 5th OT (250) & $+24$ & 4.7 & Noticeable \\
300 & 6th OT (300) & $\pm 0$ & 7.0 & \textbf{Strong -- critical} \\
325 & 6th OT (300) & $+25$ & 5.0 & Noticeable \\
400 & 8th OT (400) & $\pm 0$ & 7.0 & \textbf{Strong -- critical} \\
\bottomrule
\end{tabular}}
\caption{Amplification factor $G$ at the nearest overtone. $f_H = 274$\,Hz: $G = 4{.}7$ --- not a direct hit, but close enough for noticeable coupling.}
\end{table}

\begin{keybox}
\textbf{Optimisation rule:} Choose $f_H$ so that it lies \textbf{between two overtones}. The worst situation is $f_H = n \cdot f_\text{reed}$. For the 50\,Hz reed with $f_H = 274$\,Hz: 5th OT = 250\,Hz and 6th OT = 300\,Hz bracket $f_H$ --- no exact hit, favourable. This rule supplements the recommendation from Doc.\,0002 Ch.\,12 (avoiding resonant overtones) with the quantitative frequency-pulling aspect.
\end{keybox}

\section{Block Diagram}

\begin{center}
\fbox{\parbox{0.9\textwidth}{\small
\textbf{Signal flow of the coupled filters:}

\medskip
$\text{[Bellows pressure]} \to \text{[Valve]} \to \underbrace{H_2(f): \text{Chamber filter}}_{\text{Band-pass } f_H = 274\,\text{Hz, } Q_H = 7} \to \text{[Gap coupler]}$

$\to \underbrace{H_1(f): \text{Reed filter}}_{\text{Band-pass } f_1 = 50\,\text{Hz, } Q_1 = 50\text{--}150} \to \text{[Sound radiation]}$

\medskip
\textbf{Feedback:} $x \to S(x) \to q \to \delta p_\text{chamber} \xrightarrow{\beta = 0{.}63\,\%} F_\text{Bernoulli} \to x$

\medskip
\textbf{Closed transfer function:} $H_\text{ges}(f) = H_1(f) / (1 - \kappa_\text{eff} \cdot H_1 \cdot H_2)$
}}
\end{center}

\section{Practical Consequences}

\textbf{1.\ Tune in the chamber:} The chamber load shifts frequency by 1--5\,cents --- more than the tuning tolerance ($\pm$1--2\,cents). Never tune on the bench.

\textbf{2.\ Chamber changes require retuning:} Any change to volume, partition wall (Doc.\,0002 Ch.\,8) or valve changes $f_H$, $Q_H$ and thus the frequency shift.

\textbf{3.\ Large chambers are more stable:} $\Delta k \sim 1/V$ --- large chambers are more tolerant of geometric variations.

\textbf{4.\ Deflection acts quadratically:} $\Delta k \sim A_\text{eff}^2$ --- unevenly deflected reeds are harder to tune.

\textbf{5.\ Push and draw reeds:} Both see the same chamber impedance (Doc.\,0002 Ch.\,7, Doc.\,0003 Ch.\,3) $\to$ same shift $\to$ tremolo beat is preserved.

\textbf{6.\ Temperature:} $C_a = V/(\rho c^2)$, and $c^2 \sim T$. At $+10$°C the frequency shift decreases by $\sim$3\,\% --- small compared to the direct temperature dependence of the reed.

\section{Summary}

\begin{keybox}
\textbf{The reed and chamber form a coupled two-filter system.} The reed is a narrow-band mechanical band-pass ($f_1 = 50$\,Hz, $Q = 50$--$150$); the chamber is a broad-band acoustic band-pass ($f_H = 274$\,Hz, $Q = 5$--$10$). Coupling occurs through the gap --- nonlinearly through the Bernoulli mechanism (Doc.\,0003), but linearisable for small-signal analysis.

\textbf{Three mechanisms} shift the frequency: (1)~Static compliance load ($+0{.}3$\,cents). (2)~Resonant amplification at overtones near $f_H$ ($+1$--$5$\,cents). (3)~Nonlinear feedback to the fundamental ($+1$--$3$\,cents). \textbf{Total shift: typically 1--5\,cents.}

\textbf{The transfer function} $H_\text{ges}(f) = H_1/(1 - \kappa_\text{eff} \cdot H_1 \cdot H_2)$ is calculable. The dependencies ($V$, $A_\text{eff}$, $f_H$, $Q_H$) are quantitatively predictable. Agreement with practical findings confirms the model.

\textbf{The reed determines the frequency --- but the chamber has a say.}
\end{keybox}


\end{document}
